% Unit 081 terjemahan Bahasa Indonesia dari chapter6.tex baris 1632--1739.
% Sumber: Methods of Algebra, Volume 2, commit
% 9a5803ff77dd3257484cb177f851a73770a59dd3 (CC BY 4.0).
% stable-unit-id: o014.aljabr2.chapter6.lhs-spectral-sequence-b
% Cakupan: bukti Teorema LHS, barisan inflasi--restriksi, rumus Hopf,
% dan struktur perkalian LHS; menutup Bagian 6.9.

% segment-id: o014.li2.chapter6.lhs-spectral-sequence-b.p001
\begin{proof}
	Untuk~(i), cukup menyatakan \eqref{eqn:LHS-derived} secara konkret melalui
	barisan spektral Grothendieck dari Teorema~\ref{prop:Grothendieck-ss}.

	Sekarang kita bahas~(ii). Jika $n=1$, prasyaratnya berlaku secara hampa dan
	barisan eksak terkait tepat merupakan barisan eksak suku berderajat rendah
	dalam Teorema~\ref{prop:Grothendieck-ss}; prinsipnya ialah perhitungan tepi
	dalam Catatan~\ref{rem:edge-computing}. Berikut ini kita meninjau dan
	memperluas argumen tersebut untuk $n$ umum. Caranya ialah memanfaatkan
	suku-suku nol dalam barisan spektral.

	Cukup bahas kasus kohomologis. Syaratnya menunjukkan bahwa barisan spektral
	dalam~(i) memenuhi
	\[ 1\leq q<n\implies
	E_2^{p,q}=E_3^{p,q}=\cdots=E_\infty^{p,q}=0. \]

	Ingat bahwa arah $d_r$ adalah
	$d_r^{p,q}:E_r^{p,q}\to E_r^{p+r,q-r+1}$ dan $E_{r+1}^{p,q}$ merupakan
	kohomologi $E_r^{p,q}$ terhadap $d_r$. Semua $d$ yang keluar dari
	$E_{n+1}^{n+1,0}$ (atau masuk ke $E_{n+1}^{0,n}$) bernilai $0$. Karena itu
	diperoleh barisan eksak horizontal dalam diagram
	\[\begin{tikzcd}[row sep=small, column sep=scriptsize]
		& E_\infty^{0,n} \arrow[phantom,d,"\simeq" description,sloped]
		& & & E_\infty^{n+1,0}
		\arrow[phantom,d,"\simeq" description,sloped] & \\
		& \vdots \arrow[phantom,d,"\simeq" description,sloped]
		& & & \vdots \arrow[phantom,d,"\simeq" description,sloped] & \\
		0 \arrow[r] & E_{n+2}^{0,n} \arrow[r] &
		E_{n+1}^{0,n} \arrow[r,"{d^{0,n}_{n+1}}" inner sep=0.6em]
		\arrow[phantom,d,"\simeq" description,sloped] &
		E_{n+1}^{n+1,0} \arrow[r]
		\arrow[phantom,d,"\simeq" description,sloped] &
		E_{n+2}^{n+1,0} \arrow[r] & 0. \\
		& & \vdots \arrow[phantom,d,"\simeq" description,sloped] &
		\vdots \arrow[phantom,d,"\simeq" description,sloped] & & \\
		& \Hm^n(H,M)^{G/H} \arrow[r,"\sim"] & E_2^{0,n} & E_2^{n+1,0} &
		\Hm^{n+1}(G/H,M^H) \arrow[l,"\sim"'] &
	\end{tikzcd}\]
	Alasan adanya isomorfisme vertikal serupa: segera setelah semua $d$ yang
	masuk ke atau keluar dari $E_r^{p,q}$ bernilai $0$, terdapat isomorfisme
	$E_r^{p,q}\simeq E_{r+1}^{p,q}$. Dengan cara yang sama diperoleh
	\[ \Hm^n(G/H,M^H)\simeq E_2^{n,0}\simeq\cdots
	\simeq E_\infty^{n,0}. \]

	Selanjutnya, ingat bahwa sasaran barisan spektral
	$\Hm^{p+q}(G,M)$ mempunyai filtrasi menurun hingga
	$\cdots\supset\mathrm{F}^k\supset\mathrm{F}^{k+1}\supset\cdots$, dengan
	$E_\infty^{p,q}\simeq\gr^p\Hm^{p+q}(G,M)$. Pada tepi kuadran pertama hal
	ini memberikan
	\begin{align*}
		E_\infty^{n+1,0}&\simeq\mathrm{F}^{n+1}
		\subset\Hm^{n+1}(G,M),\\
		E_\infty^{0,n}&\simeq\Hm^n(G,M)/\mathrm{F}^1.
	\end{align*}

	Bagaimana menentukan $\mathrm{F}^1\subset\Hm^n(G,M)$? Pada garis
	$p+q=n$, hanya $(n,0)$ dan $(0,n)$ yang mungkin memenuhi
	$E_\infty^{p,q}\neq0$. Jadi filtrasi $\Hm^n(G,M)$ berbentuk
	\[ 0=\mathrm{F}^{n+1}
	\underbracket{\subset}_{E_\infty^{n,0}}\mathrm{F}^n
	=\cdots=\mathrm{F}^1
	\underbracket{\subset}_{E_\infty^{0,n}}\mathrm{F}^0
	=\Hm^n(G,M). \]
	Menggabungkan semua isomorfisme dan barisan eksak di atas memberikan hasil
	yang diminta.

	Pernyataan~(iii) sesuai dengan yang diharapkan, tetapi argumen berikut agak panjang dan
	dapat dilewati pembaca. Agar ringkas, kita hanya membahas
	$\Hm^n(G,M)\to\Hm^n(H,M)^{G/H}$; kasus lainnya serupa atau mudah. Ingat
	konstruksi dalam \S\ref{sec:double-cplx-ss}. Ambil resolusi injektif
	$M\to I$, dengan $I$ kompleks modul-$G$, lalu ambil resolusi
	Cartan--Eilenberg $J$ dari kompleks modul-$G/H$ $I^H$
	(Teorema~\ref{prop:CE-resolution}). Dengan memandang kompleks $I^G$ dan
	sebagainya sebagai bikompleks yang terkonsentrasi pada sumbu horizontal,
	diperoleh diagram komutatif bikompleks modul-$\Bbbk$
	\begin{equation}\label{eqn:LHS-SS-aux}\begin{tikzcd}
		& J^{G/H} \arrow[hookrightarrow,r] & J \\
		I^G \arrow[equal,r] & (I^H)^{G/H}
		\arrow[hookrightarrow,r] \arrow[u] & I^H. \arrow[u]
	\end{tikzcd}\end{equation}

	Mengikuti konvensi Proposisi~\ref{prop:double-cplx-ss}, operasi mengambil
	kohomologi bikompleks secara horizontal lalu vertikal ditulis
	$\Hm_{\mathrm{II}}\Hm_{\mathrm{I}}(\cdot)$. Perhatikan hal-hal berikut.
	\begin{itemize}
		\item Karena kohomologi horizontal $\Hm_{\mathrm{I}}(J)$ dari resolusi
		Cartan--Eilenberg memberikan resolusi injektif dari $\Hm^p(I^H)$ pada
		kolom ke-$p$ ($p\in\Z$), morfisme $I^H\to J$ menginduksi isomorfisme
		\[ \Hm^n(H,M)=\Hm^n(I^H)
		=\Hm_{\mathrm{II}}\Hm_{\mathrm{I}}(I^H)^{n,0}
		\rightiso\Hm_{\mathrm{II}}\Hm_{\mathrm{I}}(J)^{n,0}. \]
		\item Catatan~\ref{rem:CE-split} memberikan
		$\Hm_{\mathrm{I}}(J^{G/H})\simeq\Hm_{\mathrm{I}}(J)^{G/H}$. Karena
		$\Hm_{\mathrm{I}}(J)$ meresolusi secara vertikal setiap
		$\Hm^p(H,M)$, sedangkan $(\cdot)^{G/H}$ eksak kiri, penerapan
		$\Hm_{\mathrm{II}}\Hm_{\mathrm{I}}(\cdot)^{n,0}$ pada
		$J^{G/H}\hookrightarrow J$ bersesuaian dengan
		\[ \text{homomorfisme penyertaan}\qquad
		\Hm^n(H,M)^{G/H}\hookrightarrow\Hm^n(H,M). \]
		\item Telah ditunjukkan bahwa $I^G\hookrightarrow I^H$ pada $\Hm^n$
		(atau setara dengan itu, pada
		$\Hm_{\mathrm{II}}\Hm_{\mathrm{I}}(\cdot)^{n,0}$) menginduksi
		homomorfisme kanonik $\Hm^n(G,M)\to\Hm^n(H,M)$.
	\end{itemize}

	Filtrasikan bikompleks kuadran pertama $J^{G/H}$ secara menurun menurut
	koordinat vertikal (langkah ke-$p$ adalah bagian dengan koordinat vertikal
	$\geq p$). Barisan spektral terkait adalah $(E_r^{p,q})_{p,q}$, dengan
	$r=0,1,\ldots$. Tuliskan $\Ker(d^n)_{n,0}$ untuk proyeksi
	$\Ker(d^n)$ dari kompleks total $\tot(J^{G/H})$ pada komponen $(n,0)$.
	Deskripsi $Z_r^{p,q}$ dalam Proposisi~\ref{prop:filtered-cplx-ss}, bersama
	pengamatan di atas, memberikan
	\begin{multline}\label{eqn:LHS-SS-aux1}
		\Ker(d^n)_{n,0}\simeq Z_\infty^{0,n}
		\subset\cdots\subset Z_2^{0,n}\\
		\twoheadrightarrow E_2^{0,n}
		=\Hm_{\mathrm{II}}\Hm_{\mathrm{I}}(J^{G/H})^{n,0}
		\simeq\Hm^n(H,M)^{G/H}.
	\end{multline}

	Sekarang kita deskripsikan
	$\Hm^n(G,M)\to\Hm^n(H,M)^{G/H}$. Pilih wakil suatu unsur
	$\Hm^n(G,M)$ dalam $\Ker[I^{n,G}\to I^{n+1,G}]$, petakan melalui
	$I^G\to J^{G/H}$ ke $\Ker(d^n)_{n,0}$, lalu ikuti
	\eqref{eqn:LHS-SS-aux1} ke
	$\Hm_{\mathrm{II}}\Hm_{\mathrm{I}}(J^{G/H})^{n,0}$. Untuk selanjutnya
	mengambil citranya dalam $\Hm^n(H,M)$, terapkan
	$\Hm_{\mathrm{II}}\Hm_{\mathrm{I}}(\cdot)^{n,0}$ pada seluruh diagram
	\eqref{eqn:LHS-SS-aux} dan komposisikan sepanjang
	\begin{tikzpicture}[scale=0.5, baseline=(O)]
		\draw[->] (-1,0)--(0,0)--(0,1)--(1,1);
		\coordinate (O) at (0,0.2);
	\end{tikzpicture}.
	Jika sebagai gantinya kita mengikuti
	\begin{tikzpicture}[scale=0.5, baseline=(O)]
		\draw[->] (-1,0)--(1,0)--(1,1);
		\coordinate (O) at (0,0.2);
	\end{tikzpicture}
	dan memakai tiga pengamatan di atas, hasilnya sama dengan homomorfisme
	kanonik $\Hm^n(G,M)\to\Hm^n(H,M)$. Ini membuktikan~(iii).
\end{proof}

% segment-id: o014.li2.chapter6.lhs-spectral-sequence-b.p002
Catatan~\ref{rem:LHS-SS-mult} akan memberikan konstruksi kanonik yang mungkin
lebih alami untuk barisan spektral Lyndon--Hochschild--Serre.

% segment-id: o014.li2.chapter6.lhs-spectral-sequence-b.p003
\index{barisan eksak inflasi--restriksi (inflation-restriction sequence)}
Barisan eksak kohomologis di atas juga disebut \emph{barisan eksak
inflasi--restriksi}, sebab homomorfisme kanonik
$\Hm^n(G/H,M^H)\to\Hm^n(G,M)$ lazim disebut homomorfisme inflasi, sedangkan
$\Hm^n(G,M)\to\Hm^n(H,M)^{G/H}$ disebut homomorfisme restriksi. Morfisme yang
lebih rumit adalah
\[ \Hm^n(H,M)^{G/H}\to\Hm^{n+1}(G/H,M^H),\qquad
	\Hm_{n+1}(G/H,M_H)\to\Hm_n(H,M)_{G/H}. \]
Keduanya disebut pemetaan \emph{transgresi}, sejalan dengan pemetaan
transgresi dalam teori kelas karakteristik.
\index{transgresi (transgression)}

% segment-id: o014.li2.chapter6.lhs-spectral-sequence-b.p004
Sebagai penerapan, kita turunkan secara singkat suatu rumus Hopf untuk
menghitung $\Hm_2(G,\Z)$ dari presentasi grup $G$. Di sini kita mengambil
$\Bbbk=\Z$ dan memberi $\Z$ struktur modul-$G$ trivial.

% segment-id: o014.li2.chapter6.lhs-spectral-sequence-b.t001
\begin{corollary}[H.\ Hopf]\label{prop:Hopf-H2}
	\index[sym1]{[,]@$[\cdot,\cdot]$}
	Misalkan $G$ grup dan $G\simeq F/R$, dengan $F$ grup bebas dan
	$R\lhd F$. Tuliskan $[F,R]$ untuk subgrup dari $R$ yang dibangun oleh
	unsur-unsur $[f,r]:=frf^{-1}r^{-1}$, dengan $f\in F$ dan $r\in R$.
	Subgrup ini normal dalam $R$ dan
	$[F,R]\subset[F,F]:=F_{\mathrm{der}}$. Terdapat isomorfisme
	\[ \Hm_2(G,\Z)\simeq(R\cap[F,F])/[F,R]. \]
\end{corollary}
\begin{proof}
	Proposisi~\ref{prop:free-group-coh} memberikan $\Hm_2(F,\Z)=0$. Jadi
	barisan eksak Teorema~\ref{prop:LHS-SS} untuk $n=1$ menjadi
	\[ 0\to\Hm_2(G,\Z)\to\Hm_1(R,\Z)_G
	\to\Hm_1(F,\Z)\to\Hm_1(G,\Z)\to0. \]
	Dengan menafsirkan $\Hm_1$ melalui Contoh~\ref{eg:group-H1-further}, kita
	memperoleh
	\[ 0\to\Hm_2(G,\Z)\to\left(\frac{R}{[R,R]}\right)_G
	\to\frac{F}{[F,F]}\to\frac{G}{[G,G]}\to0. \]

	Aksi $G$ pada $R/[R,R]$ adalah sebagai berikut. Pilih prapeta $f\in F$ dari
	$g\in G$; dapat diperiksa bahwa aksi $g$ berasal dari aksi konjugasi
	$f$ pada $R$ (latihan). Untuk $\bar r\in R/[R,R]$ sembarang, pilih prapeta
	$r\in R$. Maka $g\bar r-\bar r$ (dengan operasi grup ditulis aditif)
	merupakan citra dari
	\[ frf^{-1}\cdot r^{-1}\in R
	\quad\text{(dengan operasi grup ditulis multiplikatif)}. \]
	Ketika $f$ dan $r$ bervariasi, unsur-unsur tersebut membangun $[F,R]$.
	Karena itu $(R/[R,R])_G\simeq R/[F,R]$; sisanya mudah.
\end{proof}

% segment-id: o014.li2.chapter6.lhs-spectral-sequence-b.r001
\begin{remark}[Struktur perkalian barisan spektral
Lyndon--Hochschild--Serre]\label{rem:LHS-SS-mult}
	Dalam kasus kohomologis, tuliskan barisan spektral dari modul-$G$ $M$
	sebagai $E_r^{p,q}(M)$. Operasi hasil kali cawan yang akan diperkenalkan
	dalam \S\ref{sec:cup-product}, diterapkan pada kohomologi $G/H$ dan $H$,
	melengkapi suku $E_2$ dengan struktur perkalian
	\[ \cup:E_2^{p_1,q_1}(M_1)\otimes E_2^{p_2,q_2}(M_2)
	\to E_2^{p_1+p_2,q_1+q_2}(M_1\otimes M_2), \]
	sedangkan sasaran konvergensi barisan spektral juga mempunyai hasil kali
	cawan, yang diterapkan pada grup $G$:
	\[ \cup:\Hm^{n_1}(G,M_1)\otimes\Hm^{n_2}(G,M_2)
	\to\Hm^{n_1+n_2}(G,M_1\otimes M_2). \]
	Pertanyaannya ialah apakah struktur ini dapat didefinisikan pada
	setiap halaman, dan apakah $d_r$ mempunyai sifat yang baik terhadapnya,
	misalnya aturan Leibniz
	\[ d_r(\alpha_1\cup\alpha_2)
	=d_r(\alpha_1)\cup\alpha_2
	+(-1)^{p_1+q_1}\alpha_1\cup(d_r\alpha_2),
	\qquad \alpha_i\in E_r^{p_i,q_i}(M_i). \]

	Dalam karya asli \cite[hlm.~118--119]{HS53}, Hochschild dan Serre menuliskan
	langsung suatu filtrasi kanonik pada kompleks standar ternormalisasi
	$\overline{C}(G,M)$ yang kompatibel dengan hasil kali cawan yang
	didefinisikan pada tingkat $\overline{C}(G,M)$. Dengan demikian, setiap
	halaman mewarisi perkalian yang memenuhi sifat-sifat di atas. Secara
	konkret, mereka mengambil filtrasi menurun
	$\overline{C}(G,M)=\mathrm{F}^0\supset\mathrm{F}^1\supset\cdots$. Jika
	$p>n$, suku berderajat $n$ dari kompleks $\mathrm{F}^p$ adalah $0$;
	jika tidak, suku berderajat $n$ dari $\mathrm{F}^p$ adalah
	\[ \left\{\begin{array}{r|l}
		f\in\overline{C}^n(G,M) &
		(g_1,\ldots,g_n)\ \text{memuat setidaknya }n-p+1\text{ entri dalam }H\\
		&\implies f(g_1,\ldots,g_n)=0
	\end{array}\right\}. \]
	Perinciannya diserahkan sebagai latihan bab ini; pembaca yang berminat atau
	memerlukannya dapat merujuk pada karya asli.

	Khususnya, $\overline{C}(G,\Bbbk)$ menjadi aljabar bergradasi diferensial
	terfiltrasi, dan barisan spektral $E_r^{p,q}(\Bbbk)$ mempunyai struktur
	perkalian terkait (Definisi~\ref{def:mult-ss} dan
	Proposisi~\ref{prop:dga-mult-ss}). Dari sudut pandang topologi, hal ini sama
	sekali tidak mengejutkan.
\end{remark}
